% chapters/ch01a_foundational_reference.tex
% Chapter 1A: Foundational Reference (v0) — complete

\chapter{NRMO Foundational Reference (v0)}
\label{ch:foundational}


\index{NRMO!Foundational Reference}
\index{Constitutional Layer}
\section*{Purpose}

This document preserves the primordial principles of NRMO and the
DAG-based worldview as a non-executable, read-only reference. It
exists to prevent silent drift, not to guide daily decisions.

\section*{Status}

\nrmobox{Document classification}{
\begin{tabular}{ll}
\textbf{Classification}: & Foundational Constraint \\
\textbf{Mutability}: & READ ONLY \\
\textbf{Execution Rights}: & NONE \\
\textbf{Update Policy}: & NEVER (superseded only by explicit fork)
\end{tabular}
}

% =====================================================================
\section{Primordial NRMO}
\label{sec:primordial-nrmo}

\subsection{Core Assertion}
\label{sec:primordial-core-assertion}

\nrmobox{Primordial NRMO core assertion}{
\centering
Survival is a hard constraint. Output, success, meaning, and value
are soft objectives.\\[0.4em]
NRMO does not maximise outcomes. NRMO eliminates irreversible ruin
states.
}

\subsection{Non-Ergodicity Premise}
\label{sec:primordial-non-ergodicity}

The system experiences only one trajectory. Expected value is
insufficient under single-path exposure. Any decision touching an
absorbing ruin state is invalid. This premise is non-negotiable.

\subsection{What NRMO Never Does}
\label{sec:primordial-never}

NRMO never:

\begin{itemize}
\item defines correct answers,
\item optimises happiness, success, or meaning,
\item enforces values or narratives,
\item accelerates decisions.
\end{itemize}

\noindent
NRMO answers exactly one question:

\begin{quote}
\centering
\textit{\textbf{Is this action survivable without irreversible loss?}}
\end{quote}

\subsection*{Appendix --- Primordial NRMO: Limitations and Evolution}
\label{sec:primordial-limitations}

Primordial NRMO is correct as a philosophy. Its core premise --- a
survival-first, non-ergodic decision-making stance --- is not an
opinion, but a necessary constraint for any system that must continue
operating over time.

However, when applied directly as an operational system, Primordial
NRMO exhibits critical limitations. It lacks explicit execution
boundaries, formalised governance, and protection against
absolutisation or interpretive overreach. In particular, a direct
implementation of Primordial NRMO risks:

\begin{itemize}
\item unbounded interpretation of ``correct'' action,
\item implicit personification of the decision system,
\item direct coupling of interpretation and action,
\item over-commitment driven by internally consistent logic.
\end{itemize}

For these reasons, the current NRMO architecture does not replace
Primordial NRMO, nor does it negate it. Instead, it translates the
original philosophy into a structurally safe, non-personified, and
non-ruinous operational form. The modern NRMO system exists precisely
to ensure that the original philosophy can be applied without
breaking the operator, the environment, or the system itself.

Primordial NRMO therefore occupies the role of \emph{foundational
philosophy}, while the current NRMO serves as its \emph{engineered
realisation under strict constraints}.

% =====================================================================
\section{DAG as Foundational Worldview}
\label{sec:dag-worldview}

\subsection{Why DAG Exists}
\label{sec:dag-why}

Reality contains dependencies, irreversibility, and one-way
transitions. Any system that ignores these properties lies.

\nrmobox{DAG declaration}{
\centering
DAG is not an algorithm. DAG is a declaration of reality.
}

\subsection{DAG Principles}
\label{sec:dag-principles}

\begin{itemize}
\item No cycles at the constitutional layer.
\item Dependencies must be explicit.
\item Irreversible nodes must be isolated.
\item Propagation must be local, never global.
\end{itemize}

\subsection{What DAG Is Not}
\label{sec:dag-is-not}

\begin{itemize}
\item Not an optimiser.
\item Not a planner.
\item Not a moral framework.
\item Not a prediction engine.
\end{itemize}

\noindent
DAG is a \emph{structural truth-constraint}.

% =====================================================================
\section{Relationship Between Primordial NRMO and DAG}
\label{sec:nrmo-dag-relationship}

\begin{table}[ht]
\centering
\begin{tabular}{lll}
\toprule
\textbf{Aspect} & \textbf{Primordial NRMO} & \textbf{DAG} \\
\midrule
Role & Governance & Structure \\
Authority & Absolute & None \\
Mutability & Immutable & Immutable \\
Execution & Forbidden & Forbidden \\
\bottomrule
\end{tabular}
\caption[Primordial NRMO and DAG]{Relationship table between
Primordial NRMO and DAG. Both are immutable, non-executive, and
absolute within their respective domains.}
\label{tab:nrmo-dag-relationship}
\end{table}

\noindent
They never execute. They never choose. They only invalidate.

% =====================================================================
\section{Separation From Operational Layers}
\label{sec:foundational-separation}

This document has \emph{no authority} over:

\begin{itemize}
\item Type ZERO,
\item Decision DAG,
\item SOP DAG,
\item OODA loops,
\item CI systems.
\end{itemize}

\noindent
Operational layers may evolve. This reference must not.

% =====================================================================
\section{Audit Usage Only}
\label{sec:foundational-audit}

This file may be consulted only when:

\begin{itemize}
\item behaviour feels guided,
\item safety feels like stagnation,
\item decisions feel pre-selected,
\item system drift is suspected.
\end{itemize}

\nrmobox{Audit question}{
\centering
\textit{\textbf{Is the system still only removing ruin, and nothing
else?}}
}

% =====================================================================
\section{Forbidden Uses}
\label{sec:foundational-forbidden}

The following uses are \textbf{forbidden}:

\begin{itemize}
\item justifying hesitation,
\item overriding current decisions,
\item arguing correctness,
\item imposing ideology.
\end{itemize}

\noindent
Using this document in any of these ways is a \emph{category error}.

% =====================================================================
\section{Fork Policy}
\label{sec:foundational-fork}

If reality contradicts these premises:

\begin{enumerate}
\item Fork the system.
\item Rename it.
\item Preserve this file unchanged.
\item Do not rewrite history.
\end{enumerate}

% =====================================================================
\section{Foundational Roles (Non-Executable)}
\label{sec:foundational-roles}

The following role definitions are \textbf{non-executable} and exist
solely to detect architectural drift.

\begin{description}[leftmargin=0pt,style=nextline]
\item[Type ZERO]
\textit{Role}: sovereignty enforcement and deviation suppression.\\
\textit{Never}: judgement, value assignment, optimisation.

\item[NRMO Core]
\textit{Role}: ruin veto and final YES/NO/HOLD issuance.\\
\textit{Never}: advice, optimisation, meaning attribution.

\item[DAG Layer (Operational)]
\textit{Role}: structural existence validation (cycles, dependencies,
irreversibility).\\
\textit{Never}: judgement, execution.

\item[Decision DAG]
\textit{Role}: visualisation and tracking of decision dependencies.\\
\textit{Never}: correctness determination.

\item[SOP DAG]
\textit{Role}: ordering and dependency constraints for procedures.\\
\textit{Never}: execution.

\item[SOP]
\textit{Role}: executable procedures for action.\\
\textit{Never}: judgement, structural modification.

\item[Parallel OODA]
\textit{Role}: hypothesis generation and option expansion.\\
\textit{Never}: act (execution).

\item[Passive Pattern]
\textit{Role}: continuous background monitoring and anomaly
triggering.\\
\textit{Never}: judgement, execution.

\item[HST-N]
\textit{Role}: constraint signalling based on subject tolerance and
load.\\
\textit{Never}: world evaluation, meaning attribution.
\end{description}

% =====================================================================
\section{Non-Personification Principle}
\label{sec:foundational-non-personification}

Lower layers of NRMO are \textbf{not personalities}. They are
\textbf{organs}.

\subsection{Correct Analogy}
\label{sec:non-personification-analogy}

\begin{table}[ht]
\centering
\begin{tabular}{ll}
\toprule
\textbf{Component} & \textbf{Anatomical analogue} \\
\midrule
NRMO & Governor (cerebral cortex / executive) \\
Type ZERO & Muscle \\
Parallel OODA & Vision \& vestibular system \\
SOP & Reflex arc \\
Passive Pattern & Skin \& immune system \\
\bottomrule
\end{tabular}
\caption[Anatomical analogues]{NRMO components mapped to anatomical
analogues. The mapping is illustrative; the architectural truth is
that each component is an organ, not a personality.}
\label{tab:nrmo-organs}
\end{table}

\subsection{Why Organs Must Not Become Personalities}
\label{sec:non-personification-rationale}

Giving personality to a muscle results in the following failure
modes:

\begin{itemize}
\item it moves on its own,
\item it assumes that faster is always better,
\item it ignores fatigue and damage.
\end{itemize}

\noindent
This is not strength. This is an accident.

The same applies to any lower layer. Personality implies judgement,
preference, and self-justification. Organs must never possess these
properties. Personification may be used temporarily as a teaching
metaphor; it must never be treated as an architectural truth.

Therefore, personality must exist \emph{only} at the governance
boundary, and nowhere else.

% =====================================================================
\section*{Closing Statement}
\label{sec:foundational-closing}

When everything else changes, one thing does not.

\nrmobox{Closing statement}{
\centering
\textit{NRMO survives not by knowing the right answer,}\\
\textit{but by refusing to die for the wrong one.}
}

\medskip
\noindent
This concludes the Foundational Reference (v0) --- READ ONLY. The
following chapters contain the \emph{Operational Specification}
(Read--Execute Layer) and the \emph{Foundations, Mechanism, and
Development Path} paper.
